\section{Aprendizaje No Supervisado y Diagnóstico}
\label{sec:no_supervisado}

A diferencia de la Sección \ref{sec:supervisado}, aquí no se predice un valor conocido: el objetivo es \textbf{descubrir estructura} en los datos —grupos de inversores con desempeño similar y registros que se desvían del comportamiento esperado— sin usar ninguna etiqueta.

\subsection*{Perfilado de inversores}

Para cada uno de los 24 inversores se construyó un vector resumen de su comportamiento durante las horas con generación: energía total acumulada, potencia media y máxima, desviación estándar de la potencia y porcentaje de tiempo en cero durante horas de luz (posible indicador de desconexiones o fallas).

\subsection*{Reducción de dimensión (PCA) y agrupamiento (K-Means)}

Los 24 perfiles (5 variables cada uno) se estandarizaron y se proyectaron a 2 componentes principales para visualización, reteniendo \textbf{90.0\%} de la varianza. Sobre los datos estandarizados se aplicó \textbf{K-Means} con $k=3$ clústeres.

\begin{figure}[H]
\centering
\includegraphics[width=0.8\textwidth]{fig_clustering_inversores.png}
\caption{Clustering de los 24 inversores por perfil de desempeño (proyección PCA 2D).}
\end{figure}

\begin{table}[H]
\centering
\begin{tabular}{lll}
\toprule
\textbf{Clúster} & \textbf{Inversores} & \textbf{Perfil} \\
\midrule
0 (n=14) & 08,09,10,11,12,13,14,15,17,18,20,22,23,24 & Mayor energía, menor \% en cero \\
1 (n=8)  & 01,02,03,05,06,16,19,21 & Energía y \% en cero intermedios \\
2 (n=2)  & 04,07 & Atípico: menor energía, mayor \% en cero, picos de potencia anómalos \\
\bottomrule
\end{tabular}
\caption{Clústeres de inversores por perfil de desempeño (K-Means, $k=3$)}
\end{table}

El clúster 2 —compuesto únicamente por los inversores \textbf{04} y \textbf{07}— se separa claramente del resto: son los dos inversores con mayor porcentaje de tiempo en cero durante horas de luz (36.5\% y 39.2\%, los dos valores más altos de toda la flota) y los únicos cuya potencia máxima registrada ($\sim$104--105 W) casi cuadruplica el máximo típico ($\sim$30 W) del resto de los inversores. Esta coincidencia —mismo par de inversores señalado tanto por el boxplot comparativo (Sección \ref{sec:descriptivo}) como por el clustering no supervisado— constituye evidencia consistente de que los inversores 04 y 07 son \textbf{candidatos prioritarios a revisión de mantenimiento}: combinan menor disponibilidad (más tiempo fuera de línea) con lecturas de potencia inconsistentes con el resto de la flota (posibles picos de medición espurios o fallas de sensor).

El clúster 0 agrupa a los inversores de mejor desempeño relativo (mayor energía acumulada, menor tiempo en cero), mientras que el clúster 1 ocupa una posición intermedia.

\subsection*{Detección de anomalías a nivel de registro}

Se aplicó \textbf{Isolation Forest} \cite{isolation_forest} sobre la relación hora del día -- potencia AC total, con una tasa de contaminación esperada de 2\%, para señalar intervalos de 5 minutos cuyo comportamiento se aparta de la envolvente diurna típica (posibles fallas puntuales, nubosidad extrema o errores de sensor).

\begin{figure}[H]
\centering
\includegraphics[width=0.85\textwidth]{fig_anomalias.png}
\caption{Registros anómalos (rojo) respecto al patrón diurno esperado de generación.}
\end{figure}

Se marcaron \textbf{8\,405} registros como anómalos (\textbf{1.99\%} del total de horas con luz, consistente con la tasa de contaminación configurada de 2\%).

\subsection*{Valor diagnóstico}

Este análisis complementa al modelo supervisado: mientras la Sección \ref{sec:supervisado} predice \emph{cuánta} potencia se genera, este análisis identifica \emph{qué} inversores o intervalos requieren atención de mantenimiento, sin necesidad de contar con historial de fallas etiquetado.
